\subsection{Extraneous Solutions to Radical Equations} 
Often, when we solve equations with an even root, something happens.
\begin{Example}
Solve $\sqrt{x+5}=x-1$
\begin{lpic}{}{8}
\regtextleft{0.25}{-7}{Check:};
\end{lpic}
What happened? When $x=\makeblanksmall$, the right-hand side of this equation was \makeblank[1.5in].
\end{Example}
\begin{fact}[Extraneous Solutions to Radical Equations]
If we take an \makeblank[1in] power of both sides of an equation, this can introduce an \term{extraneous solution} if one side of the equation is \makeblank[1.5in].

This happens because the inverses---the \makeblank[1in] power functions---are not \makeblank.
\end{fact}
As we could for rational equations, we can handle the possibility of extraneous solutions in two ways:
\begin{itemize}
\item Flag values that might result in extraneous solution.
\begin{itemize}
\item In this example, the only valid solutions have the right-hand side positive, so $\makeblank[1in]\geq 0$
\end{itemize}
\item Check solutions in the original equation.
\end{itemize}
For radical equations, it's usually easier to check.
